% from main.tex
%%%%%%%%%%  scan idx 398  |  folio 385  %%%%%%%%%%
il
suffit, d'après la définition de $\bar f^{*}_{U,s}$ comme adjoint à gauche du foncteur image directe, de définir un homomorphisme $f_{U,s*}(\mathcal O_X)$-linéaire
\[
  M_{(s),U}\longrightarrow f_{U,s*}\bigl(F|_{D_+(U,s)}\bigr).
\]
Pour ce faire, il faut se donner pour tout ouvert $V$ de $U$ un homomorphisme :
\begin{equation}\tag{2.5.2}
  \Gamma\bigl(V,M\bigr)_{(s)}\longrightarrow \Gamma\bigl(D_+(V,s),F\bigr)
\end{equation}
compatible aux restrictions aux ouverts de $U$, et prendre l'homomorphisme correspondant sur les faisceaux associés. L'homomorphisme (2.5.2) est défini de la façon suivante : soit $z/s^n\in\Gamma(V,M)_{(s)}$, où, $s$ étant de degré $d$, $z\in\Gamma(V,M_{nd})=\Gamma\bigl(D_+(V,s),F(nd)\bigr)$; on associe alors à $z/s^n$ la section $z.\alpha_d(s^n)^{-1}$ de $F$ sur $D_+(V,s)$, en notant $\alpha_d$ l'homomorphisme canonique
\[
  \Gamma(V,B_d)\longrightarrow \Gamma\bigl(f^{-1}(V),\mathcal O_X(d)\bigr)
\]
(cf. EGA II 2.6.3 et 2.6.4). On vérifie sans peine les compatibilités nécessaires, d'où par recollement l'homomorphisme (2.5.1).

Il est clair que dans le cas ou $F$ est quasi-cohérent l'homomorphisme ainsi construit est celui de EGA II (3.3.4.1).

L'homomorphisme de foncteur $\operatorname{Proj}_0\circ F_k\to \operatorname{Id}$ sur la catégorie des $\mathcal O_X$-Modules donne par passage aux catégories dérivées un homomorphisme de foncteur sur $D^+(\mathcal A)$ :
\begin{equation}\tag{2.5.3}
  \operatorname{Proj}_0\circ\mathbf R F_k\longrightarrow \operatorname{Id}_{D^+(\mathcal A)}.
\end{equation}

\textbf{Proposition 2.6.} Soient $S$ un schéma quasi-compact, $E$ un $\mathcal O_S$-Module localement libre de type fini, $X=\mathbf P(E)$, et $L^{\bullet}$ un complexe parfait sur $X$. Alors il existe un entier $k_0$ tel que pour $k\geq k_0$ les propriétés suivantes soient vérifiées :
\[
\begin{array}{ll}
